\documentclass[11pt]{article}
\usepackage[paperwidth=11in,paperheight=8.5in,margin=0.25in]{geometry}
\usepackage[T1]{fontenc}
\usepackage[utf8]{inputenc}
\usepackage[english]{babel}
\usepackage{lmodern}
\usepackage{microtype}
\usepackage{amsmath,amssymb}
\usepackage{xcolor}
\usepackage{tikz}
\usetikzlibrary{calc}
\hyphenpenalty=8000
\exhyphenpenalty=8000
\emergencystretch=2em
\definecolor{QCBlue}{HTML}{1F4E79}
\definecolor{QCGreen}{HTML}{2E6B5F}
\definecolor{QCOrange}{HTML}{B65B2A}
\definecolor{QCPurple}{HTML}{4B3F72}
\definecolor{QCLight}{HTML}{F4F6F8}
\definecolor{QCBorder}{HTML}{D0D6DC}
\definecolor{QCText}{HTML}{111111}
\newlength{\Margin} \setlength{\Margin}{0.25in}
\newlength{\Gap} \setlength{\Gap}{0.1in}
\newlength{\TitleH} \setlength{\TitleH}{0.8in}
\newlength{\HeaderH} \setlength{\HeaderH}{0.38in}
\newlength{\Pad} \setlength{\Pad}{0.1in}
\newlength{\UsableW}
\newlength{\UsableH}
\newlength{\CellW}
\newlength{\CellH}
\setlength{\UsableW}{\dimexpr\paperwidth-2\Margin\relax}
\setlength{\UsableH}{\dimexpr\paperheight-2\Margin\relax}
\setlength{\CellW}{5.2in}
\setlength{\CellH}{3.5in}
\newcommand{\QuadBox}[4]{%
  \node[anchor=north west, draw=QCBorder, line width=0.6pt, rounded corners=10pt,
        fill=white, minimum width=\CellW, minimum height=\CellH, inner sep=0pt] (B) at #1 {};
  \fill[#2, rounded corners=10pt] ([xshift=0pt,yshift=0pt]B.north west) rectangle ([xshift=0pt,yshift=-\HeaderH]B.north east);
  \node[anchor=west, text=white, font=\bfseries] at ([xshift=0.18in,yshift=-0.22in]B.north west) {#3};
  \node[anchor=north west, text=QCText] at ([xshift=\Pad,yshift=-\HeaderH-\Pad]B.north west) {\begin{minipage}[t]{\dimexpr\CellW-2\Pad\relax}\raggedright\small #4\end{minipage}};
}
\pagestyle{empty}
\begin{document}
\begin{tikzpicture}[remember picture,overlay]
\coordinate (NW) at ([xshift=\Margin,yshift=-\Margin]current page.north west);
\node[anchor=north west, draw=QCBorder, line width=0.6pt, rounded corners=12pt, fill=QCLight, minimum width=\UsableW, minimum height=\TitleH, inner sep=0pt] (T) at (NW) {};
\node[anchor=north west] at ([xshift=0.22in,yshift=-0.22in]T.north west) {\begin{minipage}[t]{\dimexpr\UsableW-0.44in\relax}{\Huge\bfseries \textcolor{QCBlue}{Laplace's Equation Formula Sheet}}\par\vspace{2pt}{\normalsize \textbf{Diego Juarez}\hfill \textbf{Computational Electrodynamics}\hfill \textbf{\today}}\end{minipage}};
\coordinate (GridNW) at ([yshift=-0.15in]T.south west);
\coordinate (TL) at (GridNW);
\coordinate (TR) at ([xshift=\CellW+\Gap]GridNW);
\coordinate (BL) at ([yshift=-\CellH-\Gap]GridNW);
\coordinate (BR) at ([xshift=\CellW+\Gap,yshift=-\CellH-\Gap]GridNW);
\QuadBox{(TL)}{QCBlue}{1) Harmonic Functions}{
\[
\boxed{\nabla^2\Phi=0}
\]
Charge-free electrostatic region:
\[
\rho=0,\qquad \mathbf E=-\nabla\Phi
\]
Cartesian Laplacian:
\[
\nabla^2=\partial_x^2+\partial_y^2+\partial_z^2
\]
Cylindrical radial form:
\[
\frac{1}{s}\frac{d}{ds}\left(s\frac{d\Phi}{ds}\right)=0
\]
Spherical radial form:
\[
\frac{1}{r^2}\frac{d}{dr}\left(r^2\frac{d\Phi}{dr}\right)=0
\]
1D solution:
\[
\Phi(x)=Ax+B
\]
}
\QuadBox{(TR)}{QCGreen}{2) Core Properties}{
\textbf{Mean value property:}

Value at a point equals the average over any sphere centered there.

\textbf{Maximum principle:}

No strict interior maxima or minima unless $\Phi$ is constant.

\textbf{Uniqueness:}

Dirichlet boundary data determine a unique solution.

\textbf{Conductor fact:}

In electrostatic equilibrium, $\Phi=\text{constant}$ inside a conductor.

\textbf{Separation ansatz:}
\[
\Phi(x,y)=X(x)Y(y), \qquad \Phi(r,\theta)=R(r)\Theta(\theta)
\]
}
\QuadBox{(BL)}{QCOrange}{3) Canonical Solutions}{
Parallel plates:
\[
\Phi(x)=V_0\frac{x}{d}, \qquad \mathbf E=-\frac{V_0}{d}\hat{\mathbf x}
\]
Exterior of conducting sphere:
\[
\Phi(r)=V_0\frac{R}{r}, \qquad r\ge R
\]
General spherical source-free radial solution:
\[
\Phi(r)=A+\frac{B}{r}
\]
General cylindrical radial solution:
\[
\Phi(s)=A\ln s+B
\]
Rectangle with zero side walls:
\[
\Phi(x,y)=\sum_{n=1}^{\infty} A_n \sin\left(\frac{n\pi x}{a}\right)\sinh\left(\frac{n\pi y}{a}\right)
\]
}
\QuadBox{(BR)}{QCPurple}{4) Boundary Data and Checks}{
Dirichlet:
\[
\Phi|_S=f
\]
Neumann:
\[
\frac{\partial \Phi}{\partial n}\Big|_S=g
\]
Mixed conditions are possible on different boundary segments.

\textbf{Check list}

1. Verify $\nabla^2\Phi=0$ in the domain.

2. Verify all boundary conditions.

3. Check uniqueness before trusting an image or series solution.

4. In a bounded region, interior extrema usually signal an error.
}
\end{tikzpicture}
\end{document}
